\documentclass[11pt]{article}
\usepackage[margin=0.85in]{geometry}
\usepackage{amsmath,amssymb,mathtools,booktabs,microtype,xcolor,fancyhdr,hyperref}
\definecolor{navy}{HTML}{17365D}\definecolor{teal}{HTML}{147D78}
\hypersetup{colorlinks=true,linkcolor=navy,urlcolor=teal}
\pagestyle{fancy}\setlength{\headheight}{14pt}\fancyhf{}\lhead{Extraspecial groups}\rhead{Proof and verification}\cfoot{\thepage}
\newcommand{\Sgm}{\mathcal S}
\title{\color{navy}\bfseries Extraspecial-Group Sigma--MGFs\\\large Odd exponent sectors and the $D_8/Q_8$ split}
\author{Computational verification note}\date{17 July 2026}
\begin{document}\maketitle
\begin{abstract}
We verify the proposed spectral formulas on all matrices of the two
extraspecial groups of order $3^3$ and on the faithful representations of
$D_8$ and $Q_8$.  These examples are minimal but discriminating: they detect
both the exponent-$p$ versus exponent-$p^2$ split and the plus/minus Arf type.
All 20 coefficient comparisons pass.
\end{abstract}

\section{Odd $p$: what the character does not see}
Let $E$ be extraspecial of order $p^{1+2n}$ and let $\rho$ be a faithful
irreducible representation of dimension $N=p^n$, with
$\rho(z)=\omega I_N$.  Central matrices contribute
\[
p^{-2n}\Pi_p\prod_j(1-t_j)^{-N}.
\]
In the exponent-$p$ case every noncentral element has all $p$th roots once
per block of size $p$, hence
\begin{equation}\label{eq:ep}
\Sgm_{E,\exp p}=p^{-2n}\Pi_p\prod_j(1-t_j)^{-N}
+(1-p^{-2n})\prod_j(1-t_j^p)^{-N/p}.
\end{equation}

In the exponent-$p^2$ case, $g^p=z^{\ell(gZ)}$ for a nonzero linear
functional $\ell:E/Z(E)\to\mathbb F_p$.  Set
$A_a(\mathbf t)=\prod_j(1-\omega^at_j^p)^{-N/p}$.  Counting fibers of
$\ell$ gives
\begin{equation}\label{eq:ep2}
\Sgm_{E,\exp p^2}=p^{-2n}\Pi_p\prod_j(1-t_j)^{-N}
+(p^{-1}-p^{-2n})A_0+p^{-1}\sum_{a=1}^{p-1}A_a.
\end{equation}
The coefficients in \eqref{eq:ep2} sum to one.  The last sum is the spectral
information missed by the ordinary character, which vanishes off the center.

For the direct exponent-$9$ test we use
$E=C_9\rtimes C_3=\langle x,y:y^{-1}xy=x^4\rangle$ and the matrices
\[
\rho(x)=\operatorname{diag}(\eta,\eta^4,\eta^7),\qquad
\rho(y)e_j=e_{j-1},\qquad \eta=e^{2\pi i/9}.
\]
The exponent-$3$ comparison is the $3$-dimensional Heisenberg representation.

\section{Extraspecial 2-groups}
For sign $\varepsilon=+1$ (plus type) or $-1$ (minus type), put $N=2^n$ and
\[
A_+=\prod_j(1-t_j^2)^{-N/2},\qquad
A_-=\prod_j(1+t_j^2)^{-N/2}.
\]
The quadratic form on $E/Z(E)$ has
$2^{2n-1}+\varepsilon2^{n-1}$ zeros and
$2^{2n-1}-\varepsilon2^{n-1}$ ones.  Removing the identity coset from the
first sector, restoring its two central lifts, and dividing by $|E|$ proves
\begin{align}\label{eq:e2}
\Sgm_{E^\varepsilon}={}&2^{-2n}\Pi_2\prod_j(1-t_j)^{-N}\notag\\
&+(\tfrac12+\varepsilon2^{-n-1}-2^{-2n})A_+
+(\tfrac12-\varepsilon2^{-n-1})A_-.
\end{align}
For $n=1$ the plus and minus representations are $D_8$ and $Q_8$.

\section{Direct coefficient test}
The verifier constructs every representation matrix, computes
$h_d(\operatorname{spec}g)$ from $\det(I-tg)^{-1}$, and averages
$\prod_i h_{\tau_i}$ over the group.  It separately expands
\eqref{eq:ep}, \eqref{eq:ep2}, or \eqref{eq:e2}.  Selected outcomes are:
\begin{center}
\begin{tabular}{@{}lrrr@{}}\toprule
group & $\tau$ & direct & formula\\\midrule
exponent $3$ & $(3)$&2&2\\
exponent $9$ & $(3)$&1&1\\
exponent $3$ & $(3,3)$&12&12\\
exponent $9$ & $(3,3)$&11&11\\
$D_8$ & $(2)$&1&1\\
$Q_8$ & $(2)$&0&0\\
$D_8$ & $(1,1)$&1&1\\
$Q_8$ & $(1,1)$&1&1\\
$D_8$ & $(2,2)$&3&3\\
$Q_8$ & $(2,2)$&3&3\\\bottomrule
\end{tabular}
\end{center}
The unequal rows demonstrate that the tests would fail if the spectral
sectors were collapsed.  All 20 built-in cases pass to $2\times10^{-8}$.

\section{Reproduction}
\begin{verbatim}
.venv/bin/python for_this_guy/quantum_matrix_groups/
  extraspecial_groups/verification.py
.venv/bin/marimo edit marimo_notebooks/extraspecial_groups.py
\end{verbatim}
Join each displayed command onto one shell line.
\end{document}
